\section{Discussion}

\paragraph{What the window can do.}
A reject window is cheap: it needs no model of the user, only a duration. On the public videos it did real work
with the blendshape detector, removing about three in five false triggers (10.5 to 4.1 per minute).

\paragraph{What it cannot do.}
The window sees the \emph{measured} duration of a closure, not the true one, so its effect belongs to the detector
it is paired with. The same 200\,ms window barely changed the EAR detector's rate (11.8 to 10.2), because that
detector measured the same blinks as longer. The result is also sensitive to details that are easy to overlook:
in an independent recomputation, narrowing the blendshape detector's hysteresis from 0.1 to 0.05 changed the
rate with the window from 4.2 to 3.3 per minute. A window therefore cannot be taken from published blink
durations; it has to be set for the signal, threshold and hysteresis in use. Even at its best it leaves about
four false triggers a minute, one every 15\,s, which is too many for typing on its own. It also sets a floor that
every deliberate closure must clear, and it relies on the camera resolving durations near that floor; at 10\,fps
a single frame is half the window. Finally, the public check covers four people in eight videos and was not
blind, so it describes these recordings, not blinks in general.

\paragraph{Why one session is not evidence.}
The session results are leads, not findings. They come from one session of one person, who also built the
system, knew the hypotheses and checked the responses. Closure durations were measured by the detectors, not
labelled, and the camera ran at 10--16\,fps instead of 30. The 960\,ms dot--dash boundary was chosen on the same
responses it was scored on, and none of the preregistered tests could be run. The session suggests what to test
next: whether deliberate short closures run longer than the default dot limit for most people, whether a
per-user timing calibration fixes this, and whether webcam blendshapes can separate a wink from a blink. It does
not answer any of these.
